\documentclass[aspectratio=169,10pt]{beamer}
\usepackage[UTF8,fontset=fandol]{ctex}
\usepackage{amsmath,amssymb,mathtools,booktabs}
\usepackage{tikz,pgfplots}
\pgfplotsset{compat=1.18}
\usepackage{lesson-theme}

\title{课程主题}
\subtitle{知识点}
\author{教师姓名}
\date{}

\begin{document}

% Choose one background system for the whole deck.
% Default below: TikZ background. Do not add generated background images.
% Image mode: generate cover-bg.png, content-bg.png, optional section-bg.png,
% and closing-bg.png separately from one visual brief. Delete the three \fill
% lines below and apply the correct image immediately before each frame, e.g.:
% \LessonBackground{assets/cover-bg.png}
% Never place a background with an overlay node after slide content.
\begin{frame}[plain]
  \begin{tikzpicture}[remember picture,overlay]
    \fill[LessonAccentSoft] (current page.south west) rectangle (current page.north east);
    \fill[LessonAccent!16] ([xshift=-42mm]current page.north east) circle (34mm);
    \fill[LessonAccent!28] ([xshift=-18mm,yshift=10mm]current page.south east) circle (22mm);
    \node[anchor=west,text width=.66\paperwidth] at ([xshift=13mm]current page.west) {
      {\usebeamerfont{title}\color{LessonInk}\inserttitle\par}
      \vspace{3mm}
      {\large\color{LessonAccent}\insertsubtitle\par}
      \vspace{8mm}
      {\color{LessonMuted}\insertauthor}
    };
  \end{tikzpicture}
\end{frame}
\LessonClearBackground

\begin{frame}{学习目标}
  \begin{itemize}
    \item 用一句话说明本节课要理解什么
    \item 用一句话说明本节课要能够完成什么
    \item 标明必要的前置知识
  \end{itemize}
\end{frame}

\begin{frame}{核心概念}
  \begin{columns}[T,onlytextwidth]
    \column{.54\textwidth}
      放置简洁定义、原文片段或核心结论。
    \column{.42\textwidth}
      放置精确图示、例子或对照内容。
  \end{columns}
  \[
    \LessonElement{s03-core-formula}{E=mc^2}
  \]
  \vfill
  \LessonTakeaway{写出这一页唯一需要学生记住的内容。}
\end{frame}

\begin{frame}{本节总结}
  \begin{enumerate}
    \item 核心概念
    \item 关键关系
    \item 应用或迁移
  \end{enumerate}
\end{frame}

\end{document}
